\subsection{Problem Statement}

ZK-rollups amortize proof generation and L1 data availability costs by batching transactions. In practice, rollup workloads are highly heterogeneous: transactions differ substantially in proving complexity, execution cost, and data footprint. Current sequencer batching strategies (e.g., fixed-size or fixed-interval batching) are largely heuristic-driven and do not explicitly optimize for this heterogeneity.

\textbf{Core question:}
\begin{quote}
How should a centralized ZK-rollup sequencer batch heterogeneous transactions to minimize total proving and posting cost while satisfying latency constraints?
\end{quote}

This experiment provides a systematic, simulation-based evaluation of batching policies under realistic heterogeneous workloads, with the goal of producing actionable guidance for rollup designers and operators.

\subsection{Scope and Assumptions}

To ensure feasibility and clarity, the experiment is explicitly bounded:

\begin{itemize}[leftmargin=1.5em]
    \item Centralized sequencer with parallel provers (industry-standard model).
    \item Single rollup instance; no cross-rollup interactions.
    \item Trace-driven simulation (no live deployment).
    \item Evaluation of a small, controlled batching policy space.
\end{itemize}

\textbf{Out of scope:}
\begin{itemize}[leftmargin=1.5em]
    \item Decentralized sequencing
    \item MEV extraction strategies
    \item Cryptographic proof construction details
\end{itemize}

\subsection{Fee and Cost Model}

Rollup transaction fees reflect resource usage across both Layer 2 (off-chain execution and proving) and Layer 1 (data availability and verification).

\subsubsection{Total Transaction Fee}

\begin{equation}
Tx_{fee} = L2_{fee} + L1_{fee}
\end{equation}

\subsubsection{Layer 2 Fee Component}

\begin{equation}
L2_{fee} = (\rho_{l2} + \delta) \times g
\end{equation}

\begin{itemize}[leftmargin=1.5em]
    \item $\rho_{l2}$: L2 base fee (EIP-1559-style).
    \item $\delta$: Priority tip to the sequencer.
    \item $g$: L2 gas used for execution and amortized proving cost.
\end{itemize}

\subsubsection{Layer 1 Fee Component}

\begin{equation}
L1_{fee} = (\rho_{blob} \cdot Scalar_{blob} \cdot b_{tx}) + C_{tx} + S_{tx}
\end{equation}

\begin{itemize}[leftmargin=1.5em]
    \item $\rho_{blob}$: Blob base fee (or calldata equivalent).
    \item $b_{tx}$: Bytes contributed by the transaction.
    \item $C_{tx}$: L1 gas for batch commitment.
    \item $S_{tx}$: L1 gas for ZK proof verification.
\end{itemize}

The simulator focuses on \emph{relative cost comparisons} across batching policies rather than precise fee prediction.

\subsection{Transaction Cost Profiles}

Transactions are categorized into three abstract but realistic classes:

\begin{itemize}[leftmargin=1.5em]
    \item \textbf{Type A (Light)}: Simple transfers
    \begin{itemize}
        \item Low proving cost
        \item Low DA footprint
    \end{itemize}

    \item \textbf{Type B (Medium)}: ERC-20 transfers or swaps
    \begin{itemize}
        \item Moderate proving cost
        \item Moderate DA footprint
    \end{itemize}

    \item \textbf{Type C (Heavy)}: Complex smart contract calls
    \begin{itemize}
        \item High proving cost
        \item Relatively small DA footprint
    \end{itemize}
\end{itemize}

Cost parameters are derived from public rollup documentation, benchmarks, and academic literature, with sensitivity analysis applied to test robustness.

\subsection{Batching Policy Space}

\subsubsection{Baseline (Agnostic) Policies}

\begin{itemize}[leftmargin=1.5em]
    \item \textbf{Fixed-Size ($N$)}: Batch closes after $N$ transactions.
    \item \textbf{Fixed-Interval ($T$)}: Batch closes every $T$ seconds.
\end{itemize}

These policies ignore transaction heterogeneity and serve as control baselines.

\subsubsection{Soft-Cap (Reactive) Policies}

Used in production rollups (e.g., zkSync Era):
https://github.com/matter-labs/zksync-era/blob/main/core/node/state_keeper/src/seal_criteria/conditional_sealer.rs

\begin{quote}
Close the batch when the first resource limit is reached.
\end{quote}

\begin{itemize}[leftmargin=1.5em]
    \item Prover cycle budget
    \item L2 gas limit
    \item DA size limit
    \item Maximum latency timeout
\end{itemize}

\subsubsection{Transaction-Aware (Proactive) Policies}

\begin{itemize}[leftmargin=1.5em]
    \item \textbf{Isolation / Side Queues}: Heavy transactions scheduled only when prover slack exists.
    \item \textbf{Greedy Bin Packing}: Batch construction as a 2D packing problem
    \begin{itemize}
        \item Dimension 1: Proving cost
        \item Dimension 2: DA size
    \end{itemize}

\section{Experimental Methodology}

\subsection{Study Objective and Scope}

The objective of the evaluation is to analyze how configuration choices influence ZK-Rollup performance under controlled and repeatable benchmark conditions. The study focuses on relative performance shifts rather than absolute production metrics. Accordingly, the prototype uses synthetic transfer-centric workloads, simplified state-transition execution, and parameterized proof-delay measurements. This scope enables batching, sequencing, data availability, and submission effects to be compared under consistent conditions without claiming production mainnet performance or full cryptographic prover throughput.

\begin{table}[t]
\centering
\caption{Experimental Configurations Used for Prototype Evaluation}
\label{tab:experiment_configs}
\resizebox{\linewidth}{!}{%
\begin{tabular}{@{}lccccc@{}}
\toprule
\textbf{Config} & \textbf{Batch Policy} & \textbf{Seq. Policy} & \textbf{DA Mode} & \textbf{Workload} & \textbf{Prover Mode} \\
\midrule
C1 & Fixed-size    & FCFS         & Calldata & Steady-state & Mock \\
C2 & Fixed-size    & Fee-priority & Calldata & Steady-state & Mock \\
C3 & Soft-cap      & Time-boost   & Calldata & Bursty       & Mock \\
C4 & Soft-cap      & FCFS         & Blob     & Bursty       & Mock \\
C5 & Tx-aware      & Fair BFT     & Blob     & Skewed       & Mock \\
C6 & Fixed-interval& Fee-priority & Off-chain& Steady-state & Mock \\
\bottomrule
\end{tabular}
}
\end{table}


\subsection{Experimental Factors and Regimes}

The benchmark varies the main experimental factors summarized in Table~\ref{tab:experiment_configs}: batch formation policy, sequencing policy, data availability mode, workload regime, and proof-delay mode. The workload regimes include balanced steady-state traffic, skewed synthetic transaction mixes, and bursty arrivals with latency thresholds. These regimes are designed to expose how queueing delay, committed throughput, data availability footprint, cost per transaction, and batch finalization time change as workload intensity and transaction composition vary.

\subsection{Batching Policy Space}

The benchmark considers three categories of batching policies. First, baseline agnostic policies close batches using either a fixed transaction count or a fixed time interval. These policies are simple and predictable, but they do not adapt to workload heterogeneity or changing fee conditions. Second, soft-cap policies close a batch when a resource threshold is reached, such as data availability size, proof-delay budget, gas limit, or maximum waiting time. This reflects how practical rollup systems often avoid exceeding execution, proving, or posting constraints. Third, transaction-aware policies construct batches based on transaction cost profiles, for example by limiting the proportion of heavy transactions or by treating batching as a multidimensional packing problem over proof-delay contribution and data availability footprint.

This policy space is used to evaluate whether simple fixed batching remains sufficient under homogeneous workloads or whether adaptive and transaction-aware policies become useful when proof-delay contribution and data footprint vary across transaction classes. The dynamic batching model introduced in Section~\ref{subsec:dynamic_batching_model} provides an analytical basis for reasoning about adaptive batch-size selection. In the present evaluation, it is used to interpret batching trade-offs rather than to claim a production-ready online optimizer.

\subsection{Workload Model}

Synthetic transactions are generated using a Poisson arrival process to model transfer-centric traffic under controlled arrival rates. Rather than replaying production traces, the benchmark uses synthetic workloads to isolate the effect of configuration changes. This allows the same workload distribution to be replayed across different batching, sequencing, and data availability settings.

The workload model abstracts transactions into light, medium, and heavy classes. Light transactions represent simple transfers with low data availability footprint and low proof-delay contribution. Medium transactions represent moderately larger state updates. Heavy transactions represent synthetic operations with higher data availability footprint or parameterized proof-delay cost. These classes are used to study how workload heterogeneity affects queueing delay, finalization time, committed throughput, and latency-threshold violations.

\subsection{Dynamic Batching Cost Model}
\label{subsec:dynamic_batching_model}

To reason about adaptive batch-size selection, batching is modeled as a cost-equivalent optimization problem. For a candidate batch size \(N\), the batcher estimates a penalty score \(U(N)\) that combines amortized Layer-1 cost, blob-based data availability cost, queueing delay, proof-delay contribution, and Layer-1 settlement delay:

\begin{equation}
U(N)=
\frac{
C_{\mathrm{fixed}}
+
F_{\mathrm{blob}}
\left\lceil
\frac{N b_{\mathrm{tx}}}{B_{\mathrm{blob}}}
\right\rceil
}{N}
+
\gamma
\left(
\frac{N}{2\lambda}
+
P(N)
+
T_{L1}
\right).
\label{eq:dynamic_batching}
\end{equation}

Here, \(C_{\mathrm{fixed}}\) denotes the fixed Layer-1 commit and verification cost per batch, \(b_{\mathrm{tx}}\) is the average data availability footprint per transaction, \(B_{\mathrm{blob}}\) is the blob capacity, and \(F_{\mathrm{blob}}\) is the cost associated with one blob. The variable \(\lambda\) denotes the transaction arrival rate, \(P(N)\) represents the total proof-generation delay or parameterized proof-delay for a batch of size \(N\), and \(T_{L1}\) denotes Layer-1 settlement delay. The parameter \(\gamma\) converts latency into a cost-equivalent penalty, allowing the model to represent different operator or user preferences between lower cost and lower latency.

The ceiling term captures the step-wise cost behavior of blob-based data availability. Since blob space is consumed in discrete units, the posting cost can increase abruptly when the serialized batch payload crosses a blob-capacity boundary. This makes the data availability cost curve jagged rather than perfectly linear. In the current prototype evaluation, \(P(N)\) is interpreted as parameterized mock-proof delay unless a full cryptographic prover benchmark is explicitly executed. Therefore, the model should be viewed as an analytical tool for studying batching trade-offs, not as evidence of production-grade prover optimization.

\subsection{Cost Model and Metrics}

The total transaction cost is modeled as the combination of Layer-2 processing cost and Layer-1 submission cost. The Layer-2 component captures simplified execution and parameterized proof-delay overhead, while the Layer-1 component captures calldata or blob posting cost, commitment cost, and verification-related overhead. Since the evaluation uses controlled synthetic workloads, the reported costs are interpreted as relative comparisons across configurations rather than direct production fee estimates.

The evaluation collects the following metrics:
\begin{itemize}
    \item \textbf{Submitted and Committed TPS:} the transaction input rate and the successfully committed transaction rate.
    \item \textbf{Average and Tail Latency:} average, P95, and P99 latency from transaction submission to batch finalization.
    \item \textbf{Cost and Gas per Transaction:} amortized submission and verification cost per committed transaction.
    \item \textbf{Proof-Delay and Finalization Time:} parameterized mock-proof delay and total batch finalization duration.
    \item \textbf{Data Availability Bytes and Queueing Delay:} posted data footprint and time spent waiting in the sequencer queue.
\end{itemize}

\begin{figure}[t]
    \centering
    \includegraphics[width=\linewidth]{experiment_pipeline.pdf}
    \caption{Benchmark execution pipeline used for controlled configuration sweeps.}
    \label{fig:experiment_pipeline}
\end{figure}

\subsection{Procedure and Hypotheses}

The experimental procedure consists of initializing the prototype, selecting a benchmark configuration, generating the synthetic workload, executing the benchmark run, collecting telemetry, and computing summary metrics, as shown in Fig.~\ref{fig:experiment_pipeline}. Each configuration is evaluated using the same workload-generation process so that observed differences in throughput, latency, data availability footprint, and cost can be attributed to configuration changes rather than workload variation.

The experiments are guided by three hypotheses. First, fixed-size or fixed-interval batching is expected to perform well under homogeneous workloads but degrade under heterogeneous transaction mixes because it does not adapt to transaction cost profiles. Second, soft-cap batching may improve robustness by avoiding data availability, proving, gas, or waiting-time thresholds, although it can leave some batch capacity unused. Third, transaction-aware or adaptive batching may improve latency and cost behavior when transaction classes differ significantly in data availability footprint or proof-delay contribution.
    
    \item \textbf{Bounded Mixing}: Upper bounds on heavy transaction ratios per batch.
\end{itemize}

\subsection{Experimental Design}

\subsubsection{Workload Regimes}

\begin{itemize}[leftmargin=1.5em]
    \item Balanced steady-state workload
    \item Increasing skew toward heavy transactions
    \item Bursty arrivals under strict latency SLAs
\end{itemize}

\subsubsection{Experiment Phases}

\textbf{Phase 1: Baseline Benchmarking}
\begin{itemize}
    \item Steady workload
    \item Fixed-size and fixed-interval batching
\end{itemize}

\textbf{Phase 2: Heterogeneity Sensitivity}
\begin{itemize}
    \item Gradual increase in Type C fraction
    \item Constant arrival rate
\end{itemize}

\textbf{Phase 3: Latency-Constrained Evaluation}
\begin{itemize}
    \item SLA: max waiting time = 30s
    \item Compare soft-cap vs. transaction-aware policies
\end{itemize}

\subsection{Prover Performance Model}

\begin{itemize}[leftmargin=1.5em]
    \item Proving time modeled using abstract cycle counts.
    \item Batch proving latency determined by the slowest sub-proof plus recursion overhead.
    \item Parallel provers with finite capacity and queueing effects.
\end{itemize}

\subsection{Evaluation Metrics}

\begin{center}
\begin{tabular}{ll}
\toprule
\textbf{Metric} & \textbf{Purpose} \\
\midrule
Average / P99 latency & User experience and SLA compliance \\
Throughput (tx/s) & System efficiency \\
Cost per transaction & Economic efficiency \\
Prover utilization & Resource efficiency \\
Expired transactions & SLA violations \\
\bottomrule
\end{tabular}
\end{center}

\subsection{Hypotheses}

\begin{itemize}[leftmargin=1.5em]
    \item H1: Agnostic batching becomes unstable as transaction heterogeneity increases.
    \item H2: Soft-cap policies improve robustness but waste prover capacity.
    \item H3: Transaction-aware batching strictly dominates naive batching beyond a heterogeneity threshold.
\end{itemize}

\subsection{Expected Contributions}

\begin{itemize}[leftmargin=1.5em]
    \item An open-source, reusable batching simulator for ZK-rollups.
    \item Empirical identification of failure regimes for mixed batching.
    \item Clear cost–latency Pareto frontiers for batching policy selection.
    \item Practical guidance for sequencer and prover pipeline design.
\end{itemize}
